\chapter[User Interaction Concept \textsuperscript{Novotny}]{User Interaction Concept}
 \label{chap:UserInteractionConcept}
\section{Menu Navigation}
The menu structure of the application is divided into four main parts. \\
\begin {itemize}
	\item {Main Menu}
	\item {Game Start Menu}
	\item {Options Dialog}
	\item {In-Game Menus}
\end {itemize}
The user interaction conecpt has been tailored to provide a comfortable user experience on desktop and mobile environments.
The layout of the menus has been designed to support an efficient navigation process. 

\subsection{Multidimensional User Interaction}
In order to provide an enjoyable user experience for both novices and experienced users, multiple ways of interacting with the user interface have been implemented.
\begin {itemize}
	\item {Mouse- and Touch-Based Navigation}
	\item {Keyboard-Based Navigation}
	\item {Sound-Based Navigation}
\end {itemize}
\paragraph{Keyboard-Based Navigation}
All parts of the user interface can be controlled using the keyboard exclusively. 

The following keys are used:

\begin{tabular}{ |l|l| }
  \hline
  \multicolumn{2}{|c|}{Key Map} \\
  \hline
  Arrow Keys (left / right) & pre-select item within focused list \\
  Arrow Keys (up / down) & select next / previous item in focus chain \\
  Tab / Back Tab & select next / previous item in focus chain \\
  Enter & confirm all selections and proceed (e.g. start game) \\
  Space & select pre-selected item \\
  ESC, Backspace & go back to previous menu \\
  \hline
\end{tabular}
\paragraph{Sound-Based Navigation}
All components of the application relevant to the user during a practice session can be controlled using the user´s instrument. \\
As the user might not have selected his desired input device by the time he switches to Sound-Based Navigation, notes played on any of the available devices is considered for this purpose.
Chords that the user plays on his instrument are recognized and mapped to functions of the user interface. Where applicable, the chords and their mapped functions are displayed on the user interface.\\
In order to make these functions easier to remember, certain types of chords are mapped to specific types of functionalities.
\begin {itemize}
	\item {Major chords cause a positive shift of focus or selection.}
	\item {Minor chords cause a negative shift of focus or selection.}
	\item {Diminished chords provide additional assistance to the user (e.g. displaying on-screen instructions)}
\end {itemize}
The chord \textbf{C Major} is globally reserved for confirming all settings of the current dialog and proceeding to the next step.\\
The chord \textbf{C Minor} is globally reserved for discarding all settings of the current dialog and going back to the last menu page. \\
\paragraph{Linear Navigation}
Linear navigation is used to guide a user to go from one step to another in a particular order. 
The idea behind linear navigation is that the user makes his selections in a predefined order or sequence. \\
The benefits of linear navigation include
\begin {itemize}
	\item { reducing the risk of the user omitting necessary information,}
	\item { making the process transparent and understandable for less experienced users,}
	\item { allowing keyboard- and soundbased navigation for experienced users.}
\end {itemize}
The concept of linear navigation has been implemented throughout the whole user interface. \\
The user can proceed to the next step by
\begin {itemize}
	\item {pressing associated keys on the keyboard (tab, arrows, ...),}
	\item {playing associated chords on one of the connected instrument.}
\end {itemize}.
\subsection{Main Menu} 
The Main Menu will be presented to the user as soon as the application has been initialized. 
It provides links to the Game Start Menu and the Option Menu.

\includegraphics[width=\textwidth]{main_menu_2}

Returning to the main menu is possible from every point of the user interface.
In most cases, the user will use this menu to get to the Game Start Menu. Therefore, the corresponding item is preselected in order to enable fast keyboard- or sound-based selection.
\subsection{Game Start Menu}
The purpose of the Game Start Menu is to define the game situation the user would like to use. It is used right before the game is started which is expected to happen multiple times within a session.

\includegraphics[width=\textwidth]{game_start_menu_2}
\paragraph{Main Flow}
The majority of user interactions inside the Game Start Menu consist of the following steps.
\begin {enumerate}
	\item {select the song, possibly using the search function to navigate through a wide selection of scores}
	\item {select the instrumental track withing the song}
	\item {select the section of the song to be practiced}
	\item {start the game}
\end {enumerate}
\paragraph{Preference Settings}
The following settings are expected to be changed on a less frequent basis:
\begin {itemize}
	\item {whether or not to use the training mode}
	\item {display option}
	\item {input device}
\end {itemize}
The menu layout is designed to make the main flow efficient and comfortable and to reduce the amount of necessary adjustments to the Preference Settings. \\
Linear navigation will guide the user through the items of the main flow, allowing corrections of previous selections while avoiding the need to skip settings that are rarely changed. The preference settings chain directly follows the main flow chain allowing smooth transitions between these sections. \\
The Preference Settings will be kept throughout the whole session. 
\paragraph{Search Function}
The application is designed to handle a wide selection of score files. In order to help the user to find the desired file quickly a search function has been implemented. 
\subsubsection{Focus, Pre-Selection and Selection}
The use of keyboard- and soundbased linear navigation requires the user to be aware which elements of the user interface are focused, pre-selected and selected at a certain time. Different types of highlighting are used to indicate the states of the elements.
\paragraph{Selection Process}
As mentioned earlier, the user interface supports keyboard and sound-based navigation. It allows both vertical and horizontal selection shifts. At multiple spots in the user interface, more than one component is placed in a row. Suppose a button and a horizontally aligned {\tt ListView} are placed in the same row. The horizontal arrows keys can be used to select an item within the {\tt ListView}. The user might however expect to be able to use these horizontal arrows keys to move to the other component placed on the left hand side or the right hand side of the {\tt ListView}.The example below demonstrates this issue. 

\includegraphics[width=\textwidth]{selection_issue_2}

In order to comply with user expectations two different stages of selection have been introduced:
\begin {itemize}
	\item {Pre-Selection}
	\item {Selection}
\end {itemize}
\paragraph{Pre-Selection}
Only items that are pre-selected can be selected using linear navigation. Pre-Selection can be moved using the horizontal arrows keys without changig the actual selection. This allows using the horizontal arrow keys when moving from one component within a row to another. Pre-Selected items are surrounded by a white border. 
\paragraph{Selection}
Pressing the Space key while a {\tt ListView} is focused will cause the pre-selected item to be selected. A highlight effect is used to indicate which item is currently selected.

In the example below the \textit{Rhythm} track of a song is selected. The \textit{Drum} track is preselected.

\includegraphics[width=\textwidth]{selection_process_2}
\subsection{Options Menu}
The Options Menu contains settings that affect the applications behaviour and are not directly relevant to the current session. 

Currently the options menu offers the following settings:
\begin{itemize}
	  \item{
		 whether or not preferences to keep preference settings for the next session
	  }
	  \item{
	  	whether or not audio inputs are offered as available devices in the game settings menu
	  }
	  \item{
	  	sampling rate for audio analysis
	  }
	  \item{
	  	default sound volume
	  }
\end{itemize}

\includegraphics[width=\textwidth]{options_menu_2}

Linear navigation has been implemented similarly to the game settings menu. Sound-Based navigation has not been implemented because the settings in this section are not expected to be changed during the vast majority of sessions. 

\subsection{Pause Menu}
The game can be paused at any time.\\
The pause menu has the following features:
\begin {itemize}
	\item {Adjusting the volume of the audio output}
	\item {Returning to the Game Start Menu}
\end {itemize}

\includegraphics[width=\textwidth]{pause_menu}

Sound navigation has not been implemented for the Pause Menu because a chord associated with calling the Pause Menu might coincide with notes contained in the song.

\subsection{Result Screen}
The Result Screen appears once a song has finished.

\includegraphics[width=\textwidth]{result_screen_2}

It has the following functions:
\begin {itemize}
	\item{Displaying the score the user has achieved}
	\item{Allowing the user to decide what to do next} 
\end {itemize}
Once the song has finished, the user has the following options on how to proceed:
\begin {itemize}
	\item {Watch a replay of his performance}
	\item {Try to play the same song again}
	\item {Go back to the Game Settings Menu}
\end {itemize}